\pagetitle{Appendix A / full original selected activity}{Impossible Cut + Escape the Wind + Watch and Wonder}
{\small\color{Ink}\textbf{Source:} Piter Garcia and Aastha, \emph{Invisible Invaders --- Monday, July 27, 2026}; source paragraph group labeled \emph{Activity 2}, from the original co-planned lesson. This is the complete chosen activity, not the entire Monday plan. All text below is \textbf{original, in black}; the new instruction in the body is blue. For the full camp-day context, see the \href{https://docs.google.com/document/d/1jc_dSOKXVhSoSbvfbivXlgfYOxrrJ2L6KwbR7LQZCLc/edit}{original teaching plan}.}\par\medskip
\begingroup\small\setlength{\parskip}{0.45em}
\textbf{Activity 2 - Impossible Cut + Escape the Wind + watch and wonder (9:32-10:10)}

AA gives the Impossible Cut challenge first and lets everyone, including Piter, try before offering measuring steps or hints. Compare 1 m, 10 cm, 1 cm, 5 mm, 1 mm, and 1 micrometer, then ask what classroom tools can and cannot make or verify. Show the intentionally inaccurate claim ``All microplastics are invisible.'' Place AGREE / NOT SURE / DISAGREE placards along a continuum from strongly agree to strongly disagree. Standing or moving is one route; pointing to or holding a placard while seated, asking a partner to place a marker, or passing and returning are equally valid. Name needed evidence and revise after the scale source: some microplastics are visible and others require magnification. The false sentence is a claim to test, not an answer to copy.

AA then leads Escape the Wind (see Monday teaching slides, Moving-Air Transport Model) with exactly 12 paper proxies inside one taped starting box and one finish line. Without touching the pieces, scholars choose cardboard, a clipboard, notebook, or paper fan and try to move them. Give each person one equal-length turn. Breath is optional, never required. Name the paper as a transport model, not microplastic evidence.

By 10:00, record the 12-piece starting count, how many crossed the finish line, the direction moved, where pieces stopped, and pieces that did not move. This is the transport pattern. Count all 12 pieces during cleanup. Scholars choose one vetted short video together, then speak, write, draw, point, dictate, or pass. Scholars may record their own exact Wonder Wall questions; AA or PG offers to scribe only when wanted.

Protect the 10:10-10:20 snack, movement, seated, or quiet reset.

Evidence of success: scholars explain that smaller scales may require different tools, moving air transported the paper model, and neither classroom model proves that a particle is plastic or establishes every environmental route.

Keep the Friday souvenir, scale sequence, and Wonder Wall visible through the break. Preview: ``Next Piter will help us combine our evidence and questions into one model.''

``Our cut paper and wind activity are models: each shows some things and leaves others out. A globe can show relative locations but not stars or exact travel distances. After the break, we will collect every scholar idea and question without judging it right or wrong; we can revise during the week.''

Puzzle piece: [SCIENCE + EVIDENCE + JUSTICE/ACTION] Combine Friday's two places, the scale and transport models, one short source, and youth questions in a power-aware before/during/after consensus model.

PICRAT: Interactive / Transform
\endgroup
